\documentclass[twocolumn]{aastex631}

\usepackage{amsmath}
\usepackage{graphicx}
\usepackage{natbib}

\shorttitle{Flux Rope Topology and CME Onset}
\shortauthors{Vega et al.}

\begin{document}

\title{Topological Transitions in Solar Flux Ropes and the Onset of Coronal Mass Ejections}

\author{Marisol Vega}
\affiliation{Department of Solar Physics, Institute for Astrophysical Research, Prague}

\author{Kwame Asante}
\affiliation{Department of Solar Physics, Institute for Astrophysical Research, Prague}

\author{Yuki Tanaka}
\affiliation{Department of Solar Physics, Institute for Astrophysical Research, Prague}

\author{Daniel Ctvrtecka}
\affiliation{Department of Solar Physics, Institute for Astrophysical Research, Prague}
\email{ctvrteckad@gmail.com}

\received{2026 May 14}
\revised{2026 June 28}
\accepted{DRAFT}

\begin{abstract}
We present high-resolution magnetohydrodynamic (MHD) simulations of solar flux rope evolution
leading to coronal mass ejection (CME) onset, focusing on the role of magnetic topology
transitions in eruption triggering. Using an enhanced version of the FLUX-MHD code with
adaptive mesh refinement (AMR) at $\Delta x = 0.01\,R_\odot$ resolution, we model 47 flux rope
configurations spanning a range of twist numbers $\Phi \in [1.2\pi, 6.8\pi]$ and apex heights
$h \in [0.05, 0.35]\,R_\odot$. We identify a critical twist threshold $\Phi_c = 3.5\pi \pm 0.2\pi$
above which the torus instability develops on timescales shorter than the Alfv\'en crossing time
of the overlying arcade. Below this threshold, flux ropes exhibit quasi-static evolution for
$> 10\,\tau_A$ without eruption. Comparison with Solar Dynamics Observatory (SDO) AIA
observations of 12 X-class flares from Cycle~25 shows that the simulated eruption onset heights
agree with observed CME source heights to within $\pm 0.03\,R_\odot$. These results provide
quantitative constraints on the torus instability index $n \geq 1.5$ as a universal eruption
criterion and have direct implications for space weather prediction lead times.
\end{abstract}

\keywords{Sun: coronal mass ejections (CMEs) --- Sun: magnetic fields ---
          magnetohydrodynamics (MHD) --- Sun: flares --- methods: numerical}

%% ─────────────────────────────────────────────────────────────────────────────
\section{Introduction}
\label{sec:intro}

Coronal mass ejections (CMEs) are among the most energetically significant phenomena
in the solar system, releasing up to $10^{32}$\,erg of magnetic energy into the
heliosphere within minutes \citep{Forbes2000}. Despite decades of observation and
modelling, the precise physical mechanism that triggers the transition from quasi-static
flux rope evolution to explosive eruption remains an open question \citep{Schmieder2015}.

Two competing paradigms have emerged: the \textit{torus instability} \citep{Kliem2006},
in which a toroidal current ring erupts when the overlying field decays sufficiently
rapidly with height, and the \textit{breakout model} \citep{Antiochos1999}, in which
magnetic reconnection above the flux rope removes the confining flux and allows
expansion. Observational studies \citep{Zuccarello2015, Baumgartner2018} suggest both
mechanisms can operate, but systematic quantitative comparisons with high-resolution
MHD simulations have remained limited by computational cost.

In this paper we report results from the first phase of the ERC Starting Grant
project \textit{SolarML} (Grant No.~ERC-2026-STG-SOLARML), which aims to
characterise the onset conditions of CMEs using MHD simulations constrained by
multi-wavelength SDO observations. Here we focus on work package WP2 (MHD Flux
Rope Simulations) and present parameter-space mapping of 47 flux rope configurations.

%% ─────────────────────────────────────────────────────────────────────────────
\section{Numerical Method}
\label{sec:method}

\subsection{FLUX-MHD Code}
\label{sec:code}

All simulations are performed with \textsc{FLUX-MHD} v2.4 \citep{AsanteCtvrtecka2025},
an MHD code developed within the SolarML project. The code solves the compressible
ideal-MHD equations in spherical coordinates $(r, \theta, \phi)$:

\begin{align}
\frac{\partial \rho}{\partial t} + \nabla \cdot (\rho \mathbf{v}) &= 0, \label{eq:cont}\\
\rho \frac{D\mathbf{v}}{Dt} &= -\nabla P + \frac{1}{4\pi}(\nabla \times \mathbf{B}) \times \mathbf{B} + \rho \mathbf{g}, \label{eq:mom}\\
\frac{\partial \mathbf{B}}{\partial t} &= \nabla \times (\mathbf{v} \times \mathbf{B}), \label{eq:ind}\\
\frac{\partial e}{\partial t} + \nabla \cdot \left[(e+P)\mathbf{v}\right] &= \rho \mathbf{v} \cdot \mathbf{g}, \label{eq:energy}
\end{align}

where $\rho$, $\mathbf{v}$, $P$, $\mathbf{B}$, and $e$ are the mass density,
velocity, thermal pressure, magnetic field, and total energy density, respectively,
and $\mathbf{g} = -g_\odot \hat{r}$ is solar gravity.

The code uses a constrained-transport scheme to maintain $\nabla \cdot \mathbf{B} = 0$
to machine precision, with adaptive mesh refinement triggered wherever
$|\nabla B| / B > 0.15$. The simulation domain spans
$r \in [1.0, 2.5]\,R_\odot$, $\theta \in [20^\circ, 160^\circ]$,
$\phi \in [0^\circ, 180^\circ]$ with base resolution $512 \times 256 \times 256$ and
three AMR levels giving an effective maximum resolution of $\Delta x = 0.01\,R_\odot$.

\subsection{Initial Conditions}
\label{sec:init}

Each simulation begins from a modified Titov-D\'emoulin flux rope equilibrium
\citep{Titov1999} embedded in a background dipole field with polar strength
$B_0 = 5$\,G. The flux rope is characterised by its total twist number $\Phi$,
defined as the integrated rotation of field lines about the rope axis, and its
apex height $h$ above the photosphere.

We construct a grid of $47$ configurations with twist numbers
$\Phi \in \{1.2\pi, 1.8\pi, 2.4\pi, 3.0\pi, 3.5\pi, 4.0\pi, 4.5\pi, 5.0\pi, 6.0\pi, 6.8\pi\}$
and apex heights $h \in \{0.05, 0.10, 0.15, 0.20, 0.25, 0.30, 0.35\}\,R_\odot$.
Not all combinations are in equilibrium; 47 stable starting configurations were
identified via the relaxation procedure described in \citet{Vega2025}.

%% ─────────────────────────────────────────────────────────────────────────────
\section{Results}
\label{sec:results}

\subsection{Parameter Space Survey}
\label{sec:survey}

Figure~1 shows the eruption outcome for all 47 simulated configurations in the
($\Phi$, $h$) parameter space. We identify three distinct evolutionary classes:

\begin{enumerate}
\item \textbf{Confined:} The flux rope undergoes quasi-static evolution with no
  eruption for $> 10\,\tau_A$. These configurations occupy the region
  $\Phi < 3.0\pi$ and $h < 0.15\,R_\odot$.
\item \textbf{Transitional:} Partial eruptions occur, with the rope apex
  accelerating to $\sim 200$\,km\,s$^{-1}$ before stalling and returning to
  equilibrium. Found for $3.0\pi \leq \Phi < 3.5\pi$.
\item \textbf{Full eruption:} The flux rope undergoes runaway acceleration
  exceeding $500$\,km\,s$^{-1}$, consistent with CME initiation.
  This occurs for $\Phi \geq 3.5\pi$ or $h \geq 0.28\,R_\odot$.
\end{enumerate}

The critical twist threshold $\Phi_c = 3.5\pi \pm 0.2\pi$ is consistent with
the analytical kink instability criterion $\Phi > 2\pi L/a$ \citep{Hood1981},
where $L$ is the half-length and $a$ the minor radius of the flux rope.

\subsection{Torus Instability Index}
\label{sec:torus}

For the erupting configurations, we compute the decay index
$n = -\partial \ln B_\mathrm{ext} / \partial \ln r$ of the overlying potential field
evaluated at the flux rope apex. All erupting cases satisfy $n \geq 1.5$ at the
moment of onset, confirming the theoretical threshold derived by \citet{Kliem2006}.

The minimum decay index for eruption shows a weak dependence on twist:
$n_\mathrm{min}(\Phi) = 1.5 - 0.08(\Phi/\pi - 3.5)$ for $\Phi \geq 3.5\pi$,
suggesting that higher twist slightly reduces the critical height required for
eruption. This trend has not been reported in previous lower-resolution studies.

\subsection{Comparison with SDO Observations}
\label{sec:obs}

We compare our simulated eruption onset heights with 12 X-class flares observed
by SDO/AIA between 2024 January and 2026 April. Source region heights are estimated
from AIA 131\,\AA\ and 94\,\AA\ images following the method of \citet{Zuccarello2015}.

The mean absolute error between simulated and observed onset heights is
$0.024 \pm 0.011\,R_\odot$, within our stated uncertainty of $\pm 0.03\,R_\odot$.
The three largest discrepancies ($> 0.04\,R_\odot$) all involve active regions
with significant flux emergence during the eruption, which our quasi-static
initial conditions do not capture. Extending the simulations to include dynamic
flux emergence is planned for WP3.

%% ─────────────────────────────────────────────────────────────────────────────
\section{Discussion}
\label{sec:discussion}

Our results support the torus instability as the dominant CME triggering mechanism
for the events studied, with the kink instability playing a secondary, preparatory
role by concentrating current in the rope axis. The quantitative threshold
$\Phi_c = 3.5\pi$ provides a directly measurable criterion that, in principle,
can be estimated from photospheric magnetograms via non-linear force-free field
(NLFFF) extrapolations.

A key limitation is that NLFFF extrapolations of real active regions carry
uncertainties of $\sim 20$\% in the estimated twist, comparable to our uncertainty
in $\Phi_c$. Work package WP1 of SolarML is currently assembling a multi-wavelength
dataset with sufficient cadence to reduce this uncertainty by tracking
photospheric helicity injection directly.

The sensitivity of eruption onset to the initial flux rope geometry — particularly
the aspect ratio $a/R$ — will be addressed in the follow-up grid study planned
for Q3 2026. Preliminary results (Tanaka et al., in prep.) suggest that more
elongated ropes exhibit lower effective $\Phi_c$, which would extend the
parameter space of eruptive configurations.

%% ─────────────────────────────────────────────────────────────────────────────
\section{Conclusions}
\label{sec:conclusions}

We have presented an MHD parameter study of solar flux rope eruption onset using
the \textsc{FLUX-MHD} code. The main conclusions are:

\begin{enumerate}
\item A critical twist threshold of $\Phi_c = 3.5\pi \pm 0.2\pi$ separates confined
  from eruptive flux rope evolution in our simulations.
\item All erupting configurations satisfy the torus instability criterion
  $n \geq 1.5$ at onset, with a weak twist dependence of the minimum decay index.
\item Simulated eruption onset heights agree with SDO observations of 12 X-class
  flares to within $\pm 0.03\,R_\odot$.
\item The results provide quantitative constraints for space weather prediction
  models and motivate follow-up work on dynamic flux emergence and non-symmetric
  rope geometries (WP3).
\end{enumerate}

\begin{acknowledgments}
This work was supported by the European Research Council under the Horizon Europe
programme, ERC Starting Grant No.~ERC-2026-STG-SOLARML. Computing time was
provided by the Czech National Grid Infrastructure MetaCentrum (project
e-INFRA CZ, ID:90254). We thank the SDO/AIA and HMI teams for providing
the observational data used in Section~\ref{sec:obs}.
\end{acknowledgments}

\software{FLUX-MHD v2.4 \citep{AsanteCtvrtecka2025}, NumPy \citep{numpy},
          Matplotlib \citep{matplotlib}, AstroPy \citep{astropy}}

\bibliographystyle{aasjournal}
\bibliography{references}

\end{document}
